% ============================================================================
%  APÉNDICE A
%
%  ATENCIÓN — Desde julio de 2021, A&A publica los apéndices COMO MATERIAL
%  CAMERA-READY. El editor no los tipografía, no los corrige y no cambia su
%  layout después de recibir la versión aceptada. Lo que compiles aquí es
%  literalmente lo que se publica.
%
%  Numeración: la clase convierte \section en "Appendix A", "Appendix B", ...
%  Las subsecciones quedan A.1, A.2. Referenciar con \appref{app:derivations}.
%
%  FLOATS EN APÉNDICES
%  Los floats deben quedar bajo su propio apéndice: ni antes del título, ni
%  después del título del siguiente. En apéndices cortos, un {figure*} o
%  {table*} genera una página en blanco. Receta recomendada por A&A:
%
%      \onecolumn
%      \section{Título del apéndice}
%      \begin{figure*}[ht!] ... \end{figure*}
%      texto ...
%      \FloatBarrier
%      \twocolumn
%
%  Método alternativo, solo si el apéndice contiene floats y nada de texto:
%  declarar el float primero y meter el \section DENTRO del float.
%
%  Si los floats se siguen escapando de su contexto, usar \onecolumn en todo
%  el bloque de apéndices.
% ============================================================================

\section{Derivations}
\label{app:derivations}

Texto del apéndice.
